\section{Related work}\label{sec:related}

The sequence of this paper touches four bodies of work: the machines that sort atoms and the way their traps are moved, the two published sequences of SLM holograms for rearrangement, the static solvers that compute any single frame, and the corrections for the device that displays the frames. Each is characterised below by what it assumes, what it gives and where it stops short for a moving trap.

\subsection{Machines that sort atoms and the assignment of atoms to sites}

Atom sorting was first demonstrated with moving tweezers steered by acousto-optic deflectors. Endres et al.\ assembled defect-free one-dimensional arrays of more than 50 atoms in less than 400 ms with real-time control of 100 tweezers, replacing lost atoms with surplus ones~\cite{endres}. Barredo et al.\ prepared fully loaded two-dimensional arrays of about 50 microtraps in arbitrary geometries, starting from half-filled matrices and moving atoms with a rapid sequence of moves that depends on the initial distribution~\cite{barredo}. Both machines move one trap, or one row, at a time, so that the number of moves grows with the number of atoms to be placed. Lee, Kim and Ahn framed the choice of which atom goes to which site as a combinatorial assignment problem and solved it with the Hungarian algorithm, which is fast and gives the optimal pairing for the cost they choose~\cite{lee2017}; the sequence of this paper takes that assignment as given and shares it with every baseline. Large processors keep the SLM static and add deflectors for the fast moves of a few atoms: a coherent 3000-qubit system reloads atoms continuously on top of a static array~\cite{chiu}, and a logical processor with up to 280 physical qubits moves atoms between the zones of a reconfigurable array~\cite{bluvstein}. These machines state the time and loss limits of a move but do not specify the hologram sequence; the sequences below do.

\subsection{Analytic moves by linear phases}

Kim et al.\ moved single atoms with holographic microtraps controlled by a liquid-crystal SLM and devised an analytical design of a flicker-free movement~\cite{kim2016}. In their construction a trap is a linear phase over a sub-plane of the SLM, and the linear combination $\alpha k_1x+(1-\alpha)k_2x$ of two linear phases is again a linear phase, so that the focal point sweeps smoothly between the two positions; to obtain several traps the modulation plane is divided into sub-planes, each with its own linear phase. They reported real-time transport with feedback control and a success rate of 99 percent for translations up to 10 micrometres. The construction is flicker-free by design, but it divides the SLM among the traps, so that every trap is formed by a fraction of the pixels, and the pattern is not the interference of all the traps over the whole aperture that a Gerchberg--Saxton hologram forms. It does not apply to a hologram computed by iteration, which is the case of this paper.

\subsection{Per-frame Gerchberg--Saxton with phase induction}

Kim, Kim, Lee and Ahn computed every hologram of a rearrangement movie with a modified Gerchberg--Saxton algorithm that adds zero-padding expansion of the hologram, weighted iteration feedback for reduced crosstalk, and phase induction for successive phase continuity~\cite{kimoe}. With it they formed defect-free two-dimensional atom arrays of various geometries with a loading probability of 0.98 for up to about 30 atoms, and they report a computational scalability up to order $10^3$. In the paper, each frame is generated by a fixed number of iterations toward the trap positions of that frame; the weighted feedback rescales the target by the realised amplitude at the traps relative to its mean; and the induction takes the phase of the previous frame as the initial phase of the next frame, with a condition that the change of phase between frames is small compared with one radian. The numerical experiment of their Fig.~3 uses 5 iterations per frame. They also describe a phase-change reduction (their Eqs.~4 and~5) that they did not use in their experiments. This is the published method against which the sequence is measured; we call it per-frame Gerchberg--Saxton with phase induction and run it as the paper describes it, with the weighted Gerchberg--Saxton rule of the engine (Section~\ref{sec:design}).

The induction differs from a prescription. It chooses the starting point of the iteration, so the phases of the traps start continuous; the optimization is then free to move them, and the phase of each trap at the end of a frame is whatever the iteration reaches. The passages of the paper that define the method, the induction and the weighted feedback, contain no equation that fixes where the trap phases should go along the trajectory and none that gives a number of frames for a stated intensity floor; the frame count is a parameter of the experiment.

\subsection{Linear phase interpolation}

Knottnerus et al.\ computed two endpoint holograms by weighted Gerchberg--Saxton, read the position, the amplitude and the optical phase of every tweezer from them by a Fourier transform, interpolated the three linearly, and inverted the transform at every step~\cite{lpi}. They report cycles of computation and display of 2.736 ms on commercial hardware, a success probability of 0.996 per rearrangement cycle, and, on a separate test setup with an 11 kHz camera, that a single spot made by holograms of the weighted algorithm nearly vanished several times during a move, while the interpolated spot with identical initial and final phase stayed above 70 percent. They also measured the survival of atoms against a phase slip per hologram: slips of up to several tenths of $\pi$ per step did not drastically change the survival, and slips around $\pi$ lost atoms. Their description of earlier work is that all holograms for rearrangement were calculated with the weighted algorithm, which leaves the phases of the tweezers free so that they can jump from hologram to hologram~\cite{lpi}. That statement describes holograms computed independently; the induction of Kim et al.\ is the published remedy for exactly that jump, and the measurement of this paper therefore compares the sequence with linear phase interpolation and with the induced sequence, not with independent frames. Linear phase interpolation does not optimize its intermediate frames: a frame is the inverse transform of interpolated spots, so the intensities of its traps are uneven.

\subsection{Static solvers that compute a single frame}

Every frame of a sequence is a hologram for a set of point traps, and the solvers for that problem are the building blocks. Di Leonardo et al.\ proposed an iterative algorithm for optimal holograms of arbitrary three-dimensional arrangements of traps and showed that the algorithms lead to a phase distribution that maximizes a performance quantifier expressed through the trap intensities, with a weighted form of unprecedented efficiency and uniformity~\cite{dileonardo}. Bianchi and Di Leonardo moved the weighted iteration to a graphics processor for real-time use~\cite{bianchi}. Polin et al.\ optimized holographic traps by direct search and introduced a statistical measure of their performance~\cite{polin}. For arrays of about a thousand foci, Kim et al.\ identified and removed an undesired phase rotation in the iterative Fourier transform algorithm and reached an intensity uniformity above 98 percent after three adaptive correction steps~\cite{kimol}; solvers of this family compute the arrays of a programmable simulator of 256 atoms~\cite{ebadi} and the array of about 6100 atoms~\cite{manetsch}. A different paper of the same year with a similar first-author name, in a different journal, is the source of this phase-fixed rule, and it is cited under its own key; it is not the rearrangement paper of Kim et al.\ discussed above. The frame update of the sequence proposed here is that phase-fixed rule with the phases locked.

\subsection{Camera feedback and the device}

Feedback closes the gap between the model and the instrument. Nogrette et al.\ held single atoms in arrays of up to 100 traps in arbitrary geometries~\cite{nogrette}, and Chew et al.\ reached a relative standard deviation of the tweezer intensity of 0.3 percent with a series of optimization methods driven by precise measurements~\cite{chew}; a feedback algorithm for multi-spot tweezers reports non-uniformity below 1.1 percent for more than a thousand traps~\cite{wang}. The device adds effects that the solvers do not model. A twisted-nematic liquid-crystal display shows fluctuations of its phase modulation on the order of milliseconds~\cite{lizana}, which is the time scale of a frame. Pixel crosstalk is modelled inside conjugate-gradient minimization~\cite{schroff} and inside a modified weighted algorithm~\cite{persson}. None of these treats the frames of a moving-trap sequence; this paper models the transition between two frames as a blend of their fields and does not model the temporal response of the device. The device matters for a moving trap in a way that it does not for a static one: a static array is corrected once by feedback, whereas a sequence is displayed frame after frame, and the transition between two frames is a period in which the liquid crystal is partway between two states. The sequence studied here is a way of choosing the frames, and it is independent of the correction that a feedback loop would apply to each of them.

\subsection{Gradient-based and learned holography}

Cost-function methods serve shaped targets: mixed-region amplitude freedom for atom traps~\cite{pasienski} and conjugate-gradient minimisation of the intensity error~\cite{harte} do not address the frames of a moving sequence either. Gradient methods with Wirtinger derivatives optimize any differentiable propagation model~\cite{chakravarthula}, and neural holography adds a calibrated model and a camera in the loop~\cite{peng}. Learned generators are faster than iteration: a convolutional network computes holograms of fixed cost orders of magnitude faster than iterative algorithms~\cite{eybposh}, a deep-learning pipeline produces three-dimensional holograms at video rates~\cite{shi}, and a generative network was used to create tweezer arrays for cold atoms without re-optimizing the hologram~\cite{ren}. These generators produce one frame; none of them states how a sequence of frames keeps the phase of a moving trap continuous. They are not used as baselines here, because the comparison of this paper concerns the way a sequence of frames is constructed, and any of these generators could supply a frame in place of the weighted iteration; what the sequence adds, the prescription of the trap phases along the move, applies to the frame update and not to the generator that computes it.

\subsection{Position of the proposed sequence}

Table~\ref{tab:related} compares the three sequences that move traps with SLM frames and the proposed one. The sequence proposed here takes the phase-fixed frame rule of Kim et al.~\cite{kimol}, the warm start across frames that the induction of Kim et al.~\cite{kimoe} also uses, and the interpolation idea of Knottnerus et al.~\cite{lpi}, and adds two things that none of the three has: a prescription of the trap phases along the whole move that the optimization cannot change, and a closed form that gives the number of frames for a stated shift of the pattern between frames. These are combinations of published pieces plus a bound, and the measurements of Section~\ref{sec:results} decide how much they are worth: against linear phase interpolation, which does not optimize its frames, and against per-frame Gerchberg--Saxton with phase induction, which does.

\begin{table}[t]
\centering\small
\caption{The ways to move traps with SLM frames, by how a frame is computed, what fixes the phase of a trap along the move, and whether the number of frames follows from a stated criterion. The last row is the proposed sequence.}
\label{tab:related}
\begin{tabular}{@{}p{3.1cm} p{3.0cm} p{3.9cm} p{3.9cm}@{}}
\toprule
Sequence & Frame & Phase of a moving trap & Frame count\\
\midrule
Analytic linear phases~\cite{kim2016} & linear phase on a sub-plane per trap & linear in time by construction; no Gerchberg--Saxton frame & set by the move\\
Linear phase interpolation~\cite{lpi} & inverse transform of interpolated spots, not optimized & linear between the endpoint phases & one Fourier unit of shift per hologram at most\\
Per-frame Gerchberg--Saxton with phase induction~\cite{kimoe} & fixed number of weighted iterations per frame & starts from the previous frame, then free & a parameter of the experiment\\
Flicker-bounded sequence (this paper) & weighted iterations, warm-started with the weights & locked to a prescribed smooth trajectory & derived from a bound on the shift between frames\\
\bottomrule
\end{tabular}
\end{table}
